\cvsection{Supervision and teaching}

\cvsubsection{Supervision and thesis committees}

\chronorow
  {}
  {PhD co-supervisor.}
  {\textcolor{secondary}{Leiden University (Netherlands), Leiden Institute of Advanced Computer Science
   (LIACS).}
   \begin{itemize}
     \item PhD candidate: David Quinn Alvarez.
     \item Supervisors at Leiden University: Marcello Bonsangue and Tobias Kappé.
   \end{itemize}}

\chronorow
  {22~Jun~2026}
  {Master's thesis reader (\emph{lecteur de mémoire}).}
  {\textcolor{secondary}{Université catholique de Louvain (Belgium), School of Mathematics.}
   \begin{itemize}
     \item Role: member of the thesis reading committee, assessing the written mémoire and participating in
       the thesis evaluation.
     \item Student: Piero Possidente (University of Turin; Erasmus thesis at Université catholique de
       Louvain).
   \end{itemize}}

\chronorow
  {1~Sep~2025~--\\present}
  {Member of PhD support committee (\emph{comité d’accompagnement}).}
  {\textcolor{secondary}{Université catholique de Louvain (Belgium), Research Institute in Mathematics and
   Physics.}
   \begin{itemize}
     \item Role: advise and support the PhD student and take part in the confirmation review in the second
       year of the PhD program.
     \item PhD Student: François Cayphas.
   \end{itemize}}
% \item PhD Starting date: 1 September 2025.

\chronorow
  {24~Oct~2022}
  {Co-supervision of Bachelor's thesis.}
  {\textcolor{secondary}{University of Salerno (Italy), Department of Mathematics.}
   \begin{itemize}
     \item Thesis: \emph{Universal constructions in Algebra and in Category Theory}.
     \item Student: Chiara Carratù.
   \end{itemize}}
% \item Supervisor: Luca Spada.
% \item Date of defense: 24 October 2022.

\chronorow
  {28~Sep~2021}
  {Co-supervision of Bachelor's thesis.}
  {\textcolor{secondary}{University of Salerno (Italy), Department of Mathematics.}
   \begin{itemize}
     \item Thesis: \emph{Lawvere's theories: an algebraic approach to Universal Algebra}.
     \item Student: Maria Bevilacqua.
   \end{itemize}}
% \item Supervisor: Luca Spada.
% \item Date of defense: 28 September 2021.

\cvsubsection{Teaching}

\chronorow
  {Spring 2025/26}
  {Course instructor.}
  {\textcolor{secondary}{Université catholique de Louvain (Belgium), School of Mathematics.}
   \begin{itemize}
     \item Course: \emph{Categorical Dualities in Logic}.
     \item Level: MSc in Mathematics.
     \item Tasks: lectures, office hours, preparation and assessment of oral examinations.
     \item 22.5 h.
   \end{itemize}}

\chronorow
  {Spring 2022/23}
  {Teaching assistant.}
  {\textcolor{secondary}{University of Salerno (Italy), Department of Chemistry and Biology.}
   \begin{itemize}
     \item Course: \emph{Mathematical Methods for Environmental Sciences}.
     \item Level: BSc in Environmental Sciences.
     \item Tasks: exercise classes.
     \item 20 h.
   \end{itemize}}

\chronorow
  {Fall 2022/23}
  {Teaching assistant.}
  {\textcolor{secondary}{University of Salerno (Italy), Department of Mathematics.}
   \begin{itemize}
     \item Course: \emph{Mathematical Logic}.
     \item Level: BSc in Mathematics.
     \item Tasks: exercise classes.
     \item 20 h.
   \end{itemize}}
% For students in the third year of the bachelor's degree in Mathematics.

\chronorow
  {Fall 2021/22}
  {Adjunct professor.}
  {\textcolor{secondary}{University of Salerno (Italy), Department of Chemistry and Biology.}
   \begin{itemize}
     \item Course: \emph{Mathematics 1}.
     \item Level: BSc in Environmental Sciences.
     \item Co-taught with Serafina Lapenta.
     \item Tasks: classes, office hours, preparation and evaluation of written exams, conducting oral
       examinations.
     \item 16 h.
   \end{itemize}}

\chronorow
  {Fall 2021/22}
  {Teaching assistant.}
  {\textcolor{secondary}{University of Salerno (Italy), Department of Mathematics.}
   \begin{itemize}
     \item Course: \emph{Mathematical Logic}.
     \item Level: BSc in Mathematics.
     \item Tasks: exercise classes.
     \item 20 h.
   \end{itemize}}
% \chronorow{2021/22}{``Cultore della materia''.}
% {University of Salerno (Italy), Department of Computer Science.
% Course: \emph{Discrete Mathematics}, for students in the first year of the bachelor's degree in Computer Science.
% Tasks: exam assistant (2h).}
% \chronorow{2021/22}{``Cultore della materia''.}
% {University of Salerno (Italy), Department of Chemistry and Biology.
% Course: \emph{Mathematics 1}, for students in the first year of the bachelor's degree in Environmental Sciences.
% Tasks: exam assistant (4 h).}

\chronorow
  {Fall 2019/20}
  {Teaching assistant.}
  {\textcolor{secondary}{University of Milan (Italy), Department of Mathematics.}
   \begin{itemize}
     \item Course: \emph{Algebra 2}.
     \item Level: BSc in Mathematics.
     \item Tasks: exercise classes, office hours, exam assistant.
     \item 35 h.
   \end{itemize}}
% exercise classes (28h), office hours (3h), exam assistant (4h).

\cvsubsection{Student-oriented seminars and mini-courses}
% \textit{\footnotesize I have delivered seminars and mini-courses for graduate students and young researchers on topics such as Stone duality, first-order logic, and algebras of logic. I also organized an extracurricular seminar on applications of group theory to the Rubik’s Cube, attended by about 80 students.}
% \textit{\small I have delivered various seminars and mini-courses for graduate students and young researchers, both in departments of Computer Science and Mathematics, on topics such as Stone duality, first-order logic, and algebras of logic. These activities also include an extracurricular seminar I organized as a Teaching Assistant for a course focused on group theory. On this occasion, I shared how group theory can be used to solve the Rubik’s Cube, attracting around 80 students. This experience highlights my passion for sharing the fascinating aspects of mathematics and contributing to the department’s educational program.}

\chronorow
  {25~Feb~2025}
  {Seminar for PhD students and postdoctoral researchers in Computer Science.}
  {\textcolor{secondary}{University of Birmingham (UK), School of Computer Science.}\par
   \textit{The algebras of classical first-order logic}.}

\chronorow
  {11~Feb~2025}
  {Seminar for PhD students and postdoctoral researchers in Computer Science.}
  {\textcolor{secondary}{University of Birmingham (UK), School of Computer Science.}\par
   \textit{The sequent calculus whose semantics admits the empty structure}.}

\chronorow
  {17 and 24~Oct~2024}
  {Mini-course (3h total) for PhD students and postdoctoral researchers in Computer Science.}
  {\textcolor{secondary}{University of Birmingham (UK), School of Computer Science.}\par
   \textit{Stone duality}.}

\chronorow
  {24~Mar~2022}
  {Seminar for graduate students in Mathematics.}
  {\textcolor{secondary}{New Mexico State University (USA), online.}\par
   \textit{Lawvere's algebraic theories}.}

\chronorow
  {28~May~2021}
  {Seminar for Master students in Mathematics.}
  {\textcolor{secondary}{University of Salerno (Italy).}\par
   \textit{The algebraic approach to Constraint Satisfaction Problems (CSP)}.}

\chronorow
  {5~Dec~2019}
  {Seminar for Bachelor students in Mathematics.}
  {\textcolor{secondary}{University of Milan (Italy).}\par
   \textit{The Rubik's Cube. Ideas to create a solution strategy}.}

\cvsubsection{Teacher training}

\chronorow
  {2019/20}
  {Italian teacher-training program ``24 CFU''}
  {Completed at the University of Milan. 24 ECTS (university credits), covering didactic methods and
   technologies, educational psychology and pedagogy.}
